\documentclass[10pt,a4paper]{amsart}
\usepackage[margin=19mm]{geometry}
\usepackage{amsmath,amssymb,mathtools}
\usepackage[scr=rsfs,cal=euler]{mathalfa}
\usepackage{xfrac}
\usepackage[unicode,hidelinks,bookmarks=false]{hyperref}
\newcommand{\ie}{i.e.\ }
\newcommand{\cf}{cf.\ }
\newcommand{\resp}{respectively\ }
\newcommand{\Subsection}{\subsection}
\providecommand{\nobreakdash}{\nobreak\hbox{-}}
\setlength{\emergencystretch}{2em}
\multlinegap0pt
\usepackage{bm}
\usepackage{bbm}
\usepackage{dsfont}
\usepackage{xspace}
\usepackage{cancel}
\usepackage{mathtools}
\usepackage{cleveref}
\usepackage{amsthm}
\makeatletter
\@ifundefined{theorem}{\newtheorem{theorem}{Theorem}}{}
\@ifundefined{lemma}{\newtheorem{lemma}[theorem]{Lemma}}{}
\@ifundefined{proposition}{\newtheorem{proposition}[theorem]{Proposition}}{}
\makeatother
\def\E{{\mathbb E}}
\def\P{{\mathbb P}}
\newcommand{\N}{\mathbb N}
\newcommand{\R}{\mathbb R}
\newcommand{\Sccc}{\bar{\mathcal S}}
\newcommand{\Scc}{\hat{\mathcal S}}
\newcommand{\X}{\mathcal X}
\newcommand{\reg}{\mathrm{r}}
\title[A regularized Kellerer theorem in arbitrary dimension]{A regularized Kellerer theorem in arbitrary dimension}
\author{Gudmund Pammer, Benjamin A. Robinson, Walter Schachermayer}
\date{Source: arXiv:2210.13847v3 (CC BY 4.0)}
\begin{document}
\maketitle

\noindent\emph{Source and licence.} A regularized Kellerer theorem in arbitrary dimension. Version arXiv:2210.13847v3 (CC BY 4.0), distributed under the Creative Commons Attribution 4.0 International licence, as recorded in the arXiv licence metadata for this exact version and shown on the abstract page https://arxiv.org/abs/2210.13847v3. Full rights notice: licenses/arxiv-2210.13847-CC-BY-4.0.txt.\par\smallskip

\noindent\textit{Editorial scope.} Counted unit: the Lemma whose source label is \texttt{lem:app-reg} in the article (source0.tex lines 1257--1262, source label \texttt{lem:app-reg}), delivered together with the article's own complete original proof of it (source0.tex lines 1264--1401, one \texttt{proof} environment of 138 lines). No mathematics is written, repaired or re-derived: statement and proof are the article's own text line for line. Required context retained uncounted: prop:ex-discontinuous (lines 588--590); prop:ex-regularized (lines 639--641). Every definition, display and label of those lines is retained unchanged. Omitted from this package and not redistributed: 1--587, 591--638, 642--1256, 1263--1263, 1402--1407. Nothing inside a retained excerpt is omitted or altered: every retained body is delivered exactly as the article's lines are written, so no omission receipt of any kind accompanies this package. Citations: the retained excerpts carry no citation command at all, so no bibliography entry of the article is transcribed and no reference is redistributed. Reference adaptation: none was needed and none was made; every reference inside the delivered unit resolves inside it, so no unresolved reference or citation is produced. Printed slips kept literal: no printed word, symbol, hypothesis or number of the counted statement or of its proof was altered. Layout and class: the article's own class is \texttt{amsart} with packages including natbib, amsaddr, amssymb, amsmath, mathtools, enumerate, xcolor, mathrsfs, enumitem, graphicx, subcaption, float, dsfont, bbm; the wrapper is the lane's pinned \texttt{amsart} editorial wrapper at 10pt on A4, which restates the theorem-like environments the retained text needs and adds the article's own macro definitions that the retained text needs (\texttt{\textbackslash E}, \texttt{\textbackslash N}, \texttt{\textbackslash P}, \texttt{\textbackslash R}, \texttt{\textbackslash Scc}, \texttt{\textbackslash Sccc}, \texttt{\textbackslash X}, \texttt{\textbackslash reg}), with extra wrapper packages (bm, bbm, dsfont, xspace, cancel, mathtools, cleveref). The delivered numbering is the unit's own. Assets: no retained span contains a figure environment, a graphics-inclusion command, a PDF-page inclusion, a table or a TikZ picture; no image file is copied into this package, the package declares no asset dependency and no image file is redistributed with it. Licence: version arXiv:2210.13847v3 of this article is distributed under the Creative Commons Attribution 4.0 International licence, as recorded in the arXiv licence metadata for this exact version and shown on the abstract page https://arxiv.org/abs/2210.13847v3. A full rights notice is delivered at licenses/arxiv-2210.13847-CC-BY-4.0.txt. Characters: every retained excerpt is pure ASCII, so the delivered engine, XeLaTeX with its default Latin Modern text font, reports no missing character.
\begin{proposition}\label{prop:ex-discontinuous}
	There exists a peacock $\mu$ on $\R^4$ such that there does not exist any \emph{Markov martingale} mimicking $\mu$.
\end{proposition}

\begin{proposition}\label{prop:ex-regularized}
	There exists a peacock $\mu$ on $\R^4$ such that there is no Markov martingale mimicking $\mu^\reg$, which is defined by $\mu^\reg_t \coloneqq \mu_t \ast \gamma^{t}$, for $t \in [0, 1]$.
\end{proposition}

\begin{lemma}\label{lem:app-reg}
	Let $n \in \N$ and set $t_0 = 2^{-(n+1)}$, $t_1 = 2^{-n}$. Then, in the setting of \Cref{prop:ex-regularized}, with $\Sccc^i_{t} \coloneqq \{\exists x \in S^i_{a_i(t)} \colon \; |x - Y_t| < t \}$ for $t \in [0, 1]$, $i = 1,2$, there exists a constant $C > 0$ such that
	   \begin{equation}
    	    \P[\Sccc^i_{t_0} \bigtriangleup \Sccc^i_{t_1}] \le C t_0.
	   \end{equation}
\end{lemma}

\begin{proof}
	Fix $t_0 = 2^{-(n+1)}$, $t_1 = 2^{-n}$ for some even integer $n \in \N \cup \{ 0 \}$. Then we have that $[2^{-(n+1)}, 2^{-n}] \in I_1$, and so $a_2(t)$ and $\mu_t |_{\X_2}$ are constant on this interval.    
    Subsequently, we will repeatedly employ the estimates
    \begin{equation}\label{eq:ex_regularized.estimates}
    \begin{split}
        & \E[|X_t|^4 + |N_{t^{14}}|^4] \le \frac{a_1(t)^2 + a_2(t)^2}{2} + 3 t^{28} \le C t^2,\\
        & \P[|N_{t^{14}}| \ge a] \le \frac{t^{14}}{a^2}, \quad a > 0.
	\end{split}
    \end{equation}

	Similarly as in \Cref{prop:ex-discontinuous}, we split the vector $Y$ into two parts where $Y^1_t$ denotes the first two coordinates and $Y^2_t$ the last two components of $Y_t$, for $t \in [0, 1]$. We aim to bound the probability of $Y^i$ leaving a small ball around $Y^i_{t_0}$ in the time interval $[2^{-(n +1)}, 2^{-n}]$.
	Also define the events $\Scc^i_{t} \coloneqq \{\exists x \in S^i_{a_i(t)} \colon \; |x - Y_t| < t^2\}$ for $t \in [0, 1]$, $i = 1,2$.
	
	First we compute the upper bound
    \begin{align}\label{eq:ex_regularized.ub}
        \E[|Y^i_t|^2] \le \E[|X_t^i|^2 + |N_{t^{14}}|^2] \le \frac12 a_i(t) + t^{14}.
    \end{align}
    In the case that $i = 2$ we get the lower bound
    \begin{align}
        \E[|Y_{t_0}^2|^2 \mathbbm 1_{\Scc_{t_0}^2}] & \ge  \E[|X_{t_0} + N_{t_0^{14}}|^2 \mathbbm 1_{\Scc_{t_0}^2}; |N_{t_0^{14}}| < t_0^2]\\
        &\ge
        \E[|X_{t_0} + N_{t_0^{14}} |^2\mathbbm 1_{\Scc_{t_0}^2}] - \E[|X_{t_0} + N_{t_0^{14}}|^2; |N_{t_0^{14}}| \ge t_0^2]
        \\
        &\ge \E[|X_{t_0}|^2\mathbbm 1_{\Scc_{t_0}^2}] - \E[|X_{t_0}|^2 + |N_{t_0^{14}}|^2; |N_{t_0^{14}}| \ge t_0^2],
    \end{align}
    using the independence of $X_{t_0}$ and $N_{t_0^{14}}$. We use this independence again, together with the Cauchy-Schwarz inequality and the estimates \eqref{eq:ex_regularized.estimates}, to show that
    
    \begin{align}
    	 \E[|Y_{t_0}^2|^2 \mathbbm 1_{\Scc_{t_0}^2}] &\ge \frac12 a_2(t_0) - \E[|X_{t_0}|^2 + |N_{t_0^{14}}|^2; |N_{t_0^{14}}| \ge t_0^2 ] \\
        &\ge \frac12 a_2(t_0) - \!\left(\P[|N_{t_0^{14}}| \ge t_0^2] \E[|X_{t_0}|^4 + |N_{t_0^{14}}|^4]\right)\!^\frac12
        \\
        &\ge \frac12 a_2(t_0) - C t_0^6.
    \end{align}
    
    Recalling that $a_2$ is constant on $[t_0,t_1]$, \eqref{eq:ex_regularized.ub} implies that $\E[|Y_{t_1}|^2] \leq \frac12 a_2(t_0) + t_1^{14}$.
    Since $Y$ is a martingale we have
    \begin{equation}
    	 \E[|Y_{t_1}^2 - Y^2_{t_0}|^2] = \E[|Y_{t_1}^2|^2 - |Y_{t_0}^2|^2] \le \E[|Y_{t_1}^2|^2] - \E[|Y_{t_0}^2|^2\mathbbm{1}_{\Scc_{t_0}^2}].
    \end{equation}
    Combining the upper and lower bounds above gives us
    \begin{equation}
        \label{eq:ex_regularized.claim1}
        \E[|Y_{t_1}^2 - Y^2_{t_0}|^2] \le \frac12 a_2(t_0) + t_1^{14} - \!\left(\frac12 a_2(t_0) - Ct_0^6\right)\! \le C_0 t_0^6.
    \end{equation}
    for some $C_0 > 0$.
    By Doob's maximal inequality we obtain
    \begin{equation}
        \label{eq:ex_regularized.2_Doob1}
        \P[\sup_{t \in [t_0,t_1]} |Y^2_t - Y_{t_0}^2| \ge t_0 / 2] \le 4 C_0 t^4_0.
    \end{equation}

    Next, we prove a similar bound to \eqref{eq:ex_regularized.claim1} when $i = 1$. We claim that, for some constant $C_1 > 0$,
    \begin{equation}
        \label{eq:ex_regularized.claim2}
        \E[|Y_{t_1}^1 - Y_{t_0}^1|^2 ; (\Scc^1_{t_0})^\textrm{c}] \le C_1 t_0^3.
    \end{equation}
    Since $Y$ is a martingale, we have
    \begin{equation}
	    \begin{split}
	    	\E[|Y_{t_1}^1 - Y_{t_0}^1|^2 ; (\Scc^1_{t_0})^\textrm{c}] & = \E[|Y^1_{t_1}|^2 - |Y^1_{t_0}|^2;(\Scc^1_{t_0})^\textrm{c}] \leq \E[|Y^1_{t_1}|^2 ;(\Scc^1_{t_0})^\textrm{c}]\\
	    	& \leq \E[|Y^1_{t_1}|^2 ;(\Scc^1_{t_0} \cup \Scc^2_{t_0})^\textrm{c}] + \E[|Y^1_{t_1}|^2 ;\Scc^2_{t_0}]
	    \end{split}
    \end{equation}
    We estimate each term separately, using the Cauchy-Schwarz inequality and the estimates \eqref{eq:ex_regularized.estimates} in each case. We first bound
    \begin{align}
        \E[|Y_{t_1}^1|^2 ; (\Scc^1_{t_0} \cup \Scc^2_{t_0})^\textrm{c}]
        \le
        \E[|Y_{t_1}|^4]^\frac12 \P[ (\Scc^1_{t_0} \cup \Scc^2_{t_0})^\textrm{c} ]^\frac{1}{2} 
        \le C t_1 \P[ |N_{t^{14}_0}| \ge t_0^2]^\frac12 \le C_2 t_0^{6}.
    \end{align}
    To bound the second term, we additionally apply \eqref{eq:ex_regularized.2_Doob1} to get
    \begin{align}
        \E[|Y_{t_1}^1|^2 ;\Scc^2_{t_0} ]
        &=
        \E[|Y_{t_1}^1|^2 ;\Scc^2_{t_0}, \; |Y^2_{t_1} - Y^2_{t_0}| \le t_0 / 2 ] +
        \E[|Y_{t_1}^1|^2 ;|Y^2_{t_1} - Y^2_{t_0}| \ge t_0 / 2]
        \\
        &\le
        \E[|Y_{t_1}^1|^2 ;(\Scc^1_{t_1})^\textrm{c}] + \!\left( \E[ |Y_{t_1}|^4 ] \P[ |Y_{t_1}^2 - Y_{t_0}^2| \ge t_0 / 2] \right)\!^\frac12
        \\
        &\le
        C_2 t_0^6 + \!\left( C t_0^2 \cdot 4 C_0 t_0^4 \right)\!^\frac12 \le C_3 t_0^3.
    \end{align}   
    Combining the two preceding inequalities yields \eqref{eq:ex_regularized.claim2}.
   	As before, Doob's maximal inequality implies
    \begin{equation}
        \label{eq:ex_regularized.2_Doob2}
        \P[\sup_{t \in [t_0,t_1]} |Y_t^1 - Y_{t_0}^1| \ge t_0 / 4; \; (\Scc^1_{t_0})^\textrm{c}] \le 16 C_1 t_0.    
    \end{equation}
    
    Note that, for $i = 1, 2$,
    \begin{equation}
    	\begin{split}
    		\P [ \Sccc^i_{t_0} \bigtriangleup \Sccc^i_{t_0}] & \leq \P [\Scc^i_{t_0} \bigtriangleup \Sccc^i_{t_1}] + \P[\Sccc^i_{t_0} \setminus (\Scc^i_{t_0} \cup \Sccc^i_{t_1})],
    	\end{split}
    \end{equation}
    and
    \begin{equation}
    	\P[\Sccc^i_{t_0} \setminus (\Scc^i_{t_0} \cup \Sccc^i_{t_1})] = \P[t_0^2 \leq |N_{t_0^{14}}| < t_0, \; |N_{t_1^{14}}| \geq t_1] \leq t_0^{10}.
    \end{equation}
    Hence, to prove the conclusion of the lemma, we only require bounds on $\P [\Scc^i_{t_0} \bigtriangleup \Sccc^i_{t_1}]$, for $i = 1, 2$.
   
   Now consider the events
    \begin{equation}
        A_{t_0} \coloneqq \{ |Y_{t_1}^2 - Y_{t_0}^2| \le t_0 / 4 \} \quad
        \mbox{and}\quad
        B_{t_0} \coloneqq \Scc^2_{t_0} \cap \{ |Y_{t_1}^1 - Y_{t_0}^1| \le t_0 / 4 \},
    \end{equation}
    and observe that \eqref{eq:ex_regularized.2_Doob1} and \eqref{eq:ex_regularized.2_Doob2} imply that
    \begin{equation}
        \P[A_{t_0}^\textrm{c}] \le 16C_0 t_0^4 \quad \mbox{and} \quad
        \P[\Scc^2_{t_0} \setminus B_{t_0}] \le 16 C_1 t_0.
    \end{equation}
   We calculate that
   \begin{align}
        \P[\Scc_{t_0}^1 \bigtriangleup \Sccc_{t_1}^1] & = \P[\Scc_{t_0}^1 \setminus (\Scc_{t_0}^1 \cap \Sccc_{t_1}^1)] + \P[\Sccc_{t_1}^1 \setminus (\Scc_{t_0}^1 \cap \Sccc_{t_1}^1)] \\
        & = \P[\Scc^1_{t_0}] + \P[\Sccc^1_{t_1}] - 2\P[\Scc^1_{t_0} \cap \Sccc^1_{t_1}]\\
        & \le 2(\P[\Sccc_{t_1}^1] - \P[\Scc_{t_0}^1 \cap \Sccc_{t_1}^1]) + t_0^{10},
   \end{align}
   since $\P[\Sccc^1_{t_1}] \geq \P[\Scc^1_{t_0}, |N_{t_0^{14}}| < t_0^2] \geq \P[\Scc^1_{t_0}] - t_0^{10}$.
   Note that we have the inclusions
   \begin{equation}
   		(\Scc_{t_0}^1)^\textrm{c} \cap \Sccc_{t_1}^1 \subseteq (\Scc^1_{t_0} \cup \Scc^2_{t_0})^\textrm{c} \cup (\Scc_{t_0}^2 \cap \Sccc_{t_1}^1), \quad \Scc^2_{t_0} \cap \Sccc^1_{t_1} \subseteq A_{t_0}^\textrm{c}.
   \end{equation}
   Hence
   \begin{align}
        \P[\Scc_{t_0}^1 \bigtriangleup \Sccc_{t_1}^1] \le 2\P[(\Scc^1_{t_0})^\textrm{c} \cap \Sccc^1_{t_1}] + t_0^{10}
        & \le 2 \P[(\Scc_{t_0}^1 \cup \Scc_{t_0}^2)^\textrm{c} \cup (\Scc_{t_0}^2 \cap \Sccc_{t_1}^1)] + t_0^{10}
        \\
        & \le 2 \P[|N_{t^{14}}| \ge t_0^2] + 2 \P[A_{t_0}^\textrm{c}] + t_0\\
        & \le 3 t_0^{10} + 32 C_0 t_0^4.
    \end{align}
    On the other hand, we have
    \begin{align}
        \P[\Scc_{t_0}^2 \bigtriangleup \Sccc_{t_1}^2] \le \P[\Scc^2_{t_0} \setminus B_{t_0}] + \P[|N_{t_0^{14}}| \geq t_0]\le C t_0.
    \end{align}
    Thus we have shown the claim in the case that $n$ is even. Due to symmetry of the arguments we also obtain the desired result when $n$ is odd.	
\end{proof}

\end{document}
